\section{The actual operator: definition and exact structure}\label{sec:operator}
Sections~\ref{sec:operator}--\ref{sec:boundary} use one object from the fully recovered scalar specification~\cite{archivev59,archivephysical}. Here ``physical'' means the arithmetic data actually prescribed by the research problem, as opposed to models in which signs, origins, or kernels may be freely selected. We first define that object, then establish its operator structure; subsequent sections distinguish what follows from generic norm information, exact Kloosterman representation, and the small support of the boundary.

\subsection{Parameters, coefficients, and kernel}
Let $x\to\infty$ and set
\[
I_x=(x/2,x]\cap\ZZ,\quad N=\lfloor x\rfloor-\lfloor x/2\rfloor,\quad
Q=x^{1/3},\quad H=x^{21/32},\quad U=x^{133/400}.
\]
The older source scale is $x=2X$, so $N=X+O(1)$. Let $\QQ=\{q\text{ prime}:Q<q\le2Q\}$. Fix $K>0$, put $z=(\log x)^K$, and define
\begin{align}
P(z)&=\prod_{p\le z}p,&W_z&=\prod_{p\le z}\frac p{p-1},&
C_z&=\prod_{p>z}\left(1-\frac1{(p-1)^2}\right),\nonumber\\
b_x^{(z)}(u)&=\1_{I_x}(u)C_zW_z\1_{(u+2,P(z))=1}
\prod_{\substack{p\mid u\\p>z}}\frac{p-1}{p-2},\label{eq:hybrid}\\
\beta(t)&=\frac{\Lambda(t)}{\log t}-\sum_{\substack{d\mid t\\d\le U}}\mu(d),\nonumber\\
w(u)&=\Lambda(u+2)-b_x^{(z)}(u).\label{eq:lanes}
\end{align}
Take $x$ sufficiently large that $z\ge2$. The parameter $K$ is fixed before the threshold is chosen; no assertion is made for growing $K$. The small-prime product includes $2$, and its forbidden residue is $-2$. The symbol $\mu$ in $\beta$ is the M\"obius function. Both actual coefficient vectors are real, but their signs may vary.

Fix a nonnegative real smooth function $\psi_+$ supported in $[-1,1]$ with integral $1$. Write $\e(t)=\exp(2\pi\ii t)$ and
\[
\kappa(y)=\int_{\RR}\psi_+(v)\e(-vy)\,dv,\qquad K_H(h)=\kappa(h/H).
\]
The profile need not be even, so $\kappa$ need not be real. Nevertheless, $\kappa(-y)=\overline{\kappa(y)}$ and $\kappa(0)=1$. The kernel is Schwartz, not identically zero for $|h|>H$. Reflection, $\psi_+(v)=\psi(-v)$, reconciles this convention with the older positive-sign Fourier convention.

\subsection{The complete scalar and row energy}
Set $a_q=q/(q-1)$ and define, on $\ell^2(I_x;\CC)$,
\[
A_q(u,t)=a_q\kappa((u-t)/H)
\big((q-1)\1_{q\mid u-t}-1\big)\1_{q\nmid u}\1_{q\nmid t}\1_{u\ne t},
\qquad A=\sum_{q\in\QQ}A_q.
\]
With $\ip fg=\sum\overline f g$, the target scalar is
\begin{equation}
\Cx=\ip w{A\beta}
=\sum_{q\in\QQ}q\sum_{\substack{u,t\in I_x\\u\ne t,\ q\nmid ut}}
w(u)\beta(t)K_H(u-t)\left(\1_{u\equiv t\ (q)}-\frac1{q-1}\right).
\label{eq:scalar}
\end{equation}
The outer factor $q$ is already present in the matrix definition and must not be inserted twice. The matrix is Hermitian, but its mixed scalar between two real vectors can still be complex.

The row energy under consideration is
\[
G(u)=\sum_{q\in\QQ}\sum_{t\in I_x}|A_q(u,t)|^2,\quad
D_G=\diag(G(u)),\quad Z=D_G^{-1/2}AD_G^{-1/2}.
\]
The square is taken separately at each prime; cross-prime terms are absent. Whenever every $G(u)>0$,
\begin{equation}
\Cx=\ip{D_G^{1/2}w}{Z D_G^{1/2}\beta}.\label{eq:readout}
\end{equation}
A smaller normalized operator norm therefore need not reduce the original scalar: the costs of both weighted vectors must also be counted.

\subsection{The reduced-residue correction}
\begin{proposition}[Exact positive-form decomposition]\label{prop:decomp}
Let $r_q(n)=q\1_{q\mid n}-1$ and define
\begin{align*}
T(u,t)&=\kappa((u-t)/H)\sum_{q\in\QQ}\sum_{a=1}^{q-1}\e(a(u-t)/q),\\
R(u,t)&=\kappa((u-t)/H)\sum_{q\in\QQ}\frac{r_q(u)r_q(t)}{q-1},\\
\Delta(u)&=\sum_{q\in\QQ}a_q(q-2)\1_{q\nmid u}.
\end{align*}
Then $A=T-R-\diag\Delta$ and, in positive-semidefinite order,
\[
R\ge0,\qquad0\le A+\diag\Delta\le T.
\]
\end{proposition}
\begin{proof}
On residues modulo $q$, put
\[
F_q(r,s)=q\1_{r=s}-1,\qquad
S_q(r,s)=a_q\big((q-1)\1_{r=s}-1\big)\1_{r\ne0}\1_{s\ne0}.
\]
When both residues are nonzero, exactly one is zero, or both are zero, respectively, $F_q-S_q$ equals $1/(q-1)$, $-1$, or $q-1$. Hence
\[
F_q-S_q=\frac{r_qr_q^*}{q-1}.
\]
Also, $F_q/q$ projects orthogonally onto the mean-zero subspace of all residues, while $S_q/q$ projects onto the subspace vanishing at zero and summing to zero over the units. Thus $0\le S_q\le F_q$.

Pullback along $I_x\to\ZZ/q\ZZ$ preserves this order and does not require equal residue multiplicities. For the diagonal matrix $D_v(n,n)=\e(-vn/H)$,
\[
[\kappa((u-t)/H)S_q(u,t)]_{u,t\in I_x}
=\int\psi_+(v)D_vS_qD_v^*\,dv.
\]
Integration against a nonnegative weight preserves the order. Sum over primes and remove the exact diagonal $S_q(n,n)=a_q(q-2)\1_{q\nmid n}$ to conclude.
\end{proof}
After diagonal deletion, $A$ need not be positive semidefinite. The negative directions below arise precisely at this stage. Dropping $\Delta$ during reassembly changes the problem.
